\entry{Semana del 01/12/2025} 

\section{Lunes 01/12/2025} 
Hoy vino de visita la delegación de Francia a visitar el Labo y otras instalaciones de la Facultad, con lo cual pospuse las mediciones para mañana. Una cosa importante es que hoy a la mañana al enchufar el filtro de agua para llenar el tanque no prendía (parece finalmente que se quemó algún componente en el motor). 

Con el agua estante que hay en el tanque alcanza para un par de mediciones en la cuba chica. 

Cuando se pueda medir en la cuba grande lo que me gustaría empezar haciendo es forzado con la paleta que abarca todo el ancho en una punta, y con los sensores ver si las ondas (que en principio deberían ser colineales) interactúan a tres o cuatro ondas, ya que teóricamente las primeras resonancias para gravedad y capilaridad deberían ser con cinco. 

Después la idea sería forzar dos ondas que interactúen entre sí en el rango capilar y poner los sensores en la dirección de la onda resultante, viendo en principio que la amplitud de la tercera aumenta en función de la distancia (y por tanto del tiempo que pasan interactuando las ondas), hasta saturar supongo. En el espectro temporal se deberían ver los picos de las ondas interactuante sy la resultante (este último pico aumentando cuanto más lejos el sensor). 

Otra alternativa es poner una sola onda en un motor y sumarle otra con otra frecuencia con una amplitud que crezca lentamente y ver cómo se van prendiendo las interacciones mediante las correlaciones que veníamos calculando.  

Además de poner los sensores no colineales con un forzado de ruido como para intentar reobtener la relación de dispersión mediante beamforming. 

\section{Martes 02/12/2025} 
Hoy estuve midiendo en la cuba chica. La idea sería sistematizar un poco más las mediciones originales (variando banda de frecuencia y amplitud), pero tomando más puntos. Además quise probar las mediciones de ruido moduladas por una exponencial que crece en amplitud hasta saturar lentamente. 

El setup que usé fue básicamente el mismo que el del día 22/08/2025, con ese mismo sensor original. La única cosa que cambió fue que usé componentes electrónicos nuevos en lugar de esos mismos. 

En base a pruebas que estuve haciendo con las primeras mediciones noté que el espectro no mejoraba mucho más tomando señales más largas de un minuto, con lo cual para poder tomar la mayor cantidad de puntos posibles esta vez tomé señales de una duración total de 2.5 minutos (1.5 minutos usables sin el transitorio). 







